% Part: incompleteness
% Chapter: arithmetization-syntax
% Section: proofs-in-lk

\documentclass[../../../include/open-logic-section]{subfiles}

\begin{document}

\olfileid{inc}{art}{plk}
\olsection{$\Log{LK}$માં \usetoken{P}{derivation}}

\begin{explain}
!!{derivation}sનું અંકગણિતીકરણ કરવા માટે
આપણે !!{derivation}sને સંખ્યાઓ વડે નિરૂપવાં પડશે.
દરેક !!{derivation} એ સિક્વન્ટોનું વૃક્ષ છે
અને તેમાં દરેક અનુમાનને એક ચિહ્નન પણ
આપેલું છે. તેથી પુનરાવર્તી નિરૂપણ
સૌથી સ્વાભાવિક છે: !!{derivation}ને
એવી ટપલ તરીકે નિરૂપીએ જેના ઘટકો
અંતિમ સિક્વન્ટ, ચિહ્નન અને છેલ્લા અનુમાનના
આધાર-સિક્વન્ટો સુધી પહોંચતી
ઉપ-!!{derivation}sનાં નિરૂપણ છે.
\end{explain}

\begin{defn}
$\Gamma$ એ !!{sentence}sની સાન્ત શ્રેણી,
$\Gamma = \tuple{!A_1, \dots, !A_n}$ હોય,
તો $\Gn{\Gamma} = \tuple{\Gn{!A_1},
  \dots, \Gn{!A_n}}$.

$\Gamma \Sequent \Delta$ સિક્વન્ટ હોય,
તો $\Gamma \Sequent \Delta$ની
ગડેલ-સંખ્યા આ છે:
\[
\Gn{\Gamma \Sequent \Delta} = \tuple{\Gn{\Gamma}, \Gn{\Delta}}
\]

$\pi$ એ $\Log{LK}$માં !!a{derivation} હોય,
તો $\Gn{\pi}$ને આ પ્રમાણે વ્યાખ્યાયિત કરીએ:
\begin{enumerate}
\item $\pi$ માત્ર આરંભિક સિક્વન્ટ
  $\Gamma \Sequent \Delta$નું બનેલું હોય,
  તો $\Gn{\pi}$ આ છે:
  \[
    \tuple{0, \Gn{\Gamma \Sequent \Delta}}.
  \]
\item $\pi$નો અંત એક કે બે આધાર-સિક્વન્ટવાળા
  અનુમાનથી થતો હોય, તેનું ફલિત-સિક્વન્ટ
  $\Gamma \Sequent \Delta$ હોય અને
  $\pi_1$ તથા, લાગુ પડે ત્યારે, $\pi_2$
  છેલ્લા અનુમાનના આધાર-સિક્વન્ટોમાં
  પૂરાં થતાં સીધાં ઉપપ્રમાણો હોય,
  તો અનુક્રમે $\Gn{\pi}$ આ છે:
  \begin{align*}
    & \tuple{1, \Gn{\pi_1}, \Gn{\Gamma \Sequent \Delta}, k} \text{ અથવા}\\ 
    & \tuple{2, \Gn{\pi_1}, \Gn{\pi_2}, \Gn{\Gamma \Sequent \Delta}, k}, 
  \end{align*}
  અહીં છેલ્લા અનુમાનમાં વપરાયેલા નિયમ
  પ્રમાણે નીચેના કોષ્ટકમાંથી~$k$ મળે છે:

  \begin{tabular}{lccccccc}
    \text{નિયમ:} & \LeftR{\Weakening} & \RightR{\Weakening} &
    \LeftR{\Contraction} & \RightR{\Contraction} &
    \LeftR{\Exchange} & \RightR{\Exchange} \\
    $k$: & 1 & 2 & 3 & 4 & 5 & 6 \\[2ex]
    \text{નિયમ:} & \LeftR{\lnot} & \RightR{\lnot} &
    \LeftR{\land} & \RightR{\land} & 
    \LeftR{\lor} & \RightR{\lor} \\
    $k$: & 7 & 8 & 9 & 10 & 11 & 12 \\[2ex]
    \text{નિયમ:} & \LeftR{\lif} & \RightR{\lif} &
    \LeftR{\lforall} & \RightR{\lforall} &
    \LeftR{\lexists} & \RightR{\lexists} \\
    $k$: & 13 & 14 & 15 & 16 & 17 & 18 \\[2ex]
    \text{નિયમ:} & \Cut & = \\
    $k$: & 19 & 20
  \end{tabular}
\end{enumerate}
\end{defn}

\begin{ex}
  આ અત્યંત સરળ !!{derivation} જુઓ:
  \begin{prooftree}
    \Axiom$!A \fCenter !A$
    \RightLabel{\LeftR{\land}}
    \UnaryInf$!A \land !B \fCenter !A$
    \RightLabel{\RightR{\lif}}
    \UnaryInf$\fCenter (!A \land !B) \lif !A$
  \end{prooftree}
  આરંભિક સિક્વન્ટની ગડેલ-સંખ્યા
  $p_0 = \tuple{0,
    \Gn{!A \Sequent !A}}$ થાય.
  $\LeftR{\land}$ના ફલિત-સિક્વન્ટમાં
  પૂરી થતી !!{derivation}ની ગડેલ-સંખ્યા
  $p_1 =
    \tuple{1, p_0, \Gn{!A \land !B \Sequent !A}, 9}$
  થાય ($1$ એટલા માટે કે
  $\LeftR{\land}$નો એક આધાર-સિક્વન્ટ છે;
  $!A \land !B \Sequent !A$ના
  ફલિત-સિક્વન્ટની ગડેલ-સંખ્યા લેવાઈ છે;
  અને $9$ એ $\LeftR{\land}$નો સંકેત
  નંબર છે). આખી !!{derivation}ની
  ગડેલ-સંખ્યા તેથી
  $\tuple{1, p_1, \Gn{\Sequent (!A \land !B) \lif !A}, 14}$,
  એટલે કે
  \[
  \tuple{1, \tuple{1, \tuple{0, \Gn{!A \Sequent !A}}, \Gn{!A \land !B \Sequent !A}, 9},
    \Gn{\Sequent (!A \land !B) \lif !A}, 14}.
  \]
  \sourcecorrection{OLINC-006}{સ્થિર અંગ્રેજી મૂળમાં
  આ ઉદાહરણની ટપલમાં બે જગ્યાએ
  સિક્વન્ટની અંદર વધારાનું બંધ કૌંસ છે.
  અહીં એ વધારાનાં કૌંસ કાઢ્યાં છે;
  નિયમ-સંકેત~$14$ અને ટપલના બાકીના
  ઘટકો યથાવત્ છે.}
\end{ex}

\begin{explain}
!!{derivation}sનું નિરૂપણ નક્કી કર્યા પછી,
આવા !!{derivation}s પર આદિમ પુનરાવર્તી રીતે ક્રિયા
કરી શકાય છે અને તેમના આવશ્યક ગુણધર્મો
તથા સંબંધો એ જ રીતે વ્યક્ત કરી શકાય છે
તે પણ બતાવવું પડશે. કેટલીક ક્રિયાઓ
સરળ છે. દાખલા તરીકે, !!a{derivation}ની
ગડેલ-સંખ્યા~$p$ હોય, તો
$\fn{EndSequent}(p) = (p)_{(p)_0+1}$
તેના અંતિમ સિક્વન્ટની ગડેલ-સંખ્યા આપે છે.
આરંભિક સિક્વન્ટ સિવાયના કિસ્સામાં
$\fn{LastRule}(p) = (p)_{(p)_0+2}$
છેલ્લા નિયમનો સંકેત આપે છે.
આરંભિક સિક્વન્ટના કિસ્સામાં આ વિધેયને
ઇચ્છીએ તો શૂન્ય મૂલ્ય આપીને સર્વત્ર
વ્યાખ્યાયિત કરી શકીએ.
\sourcecorrection{OLINC-007}{સ્થિર અંગ્રેજી મૂળ
અહીં અંતિમ સિક્વન્ટના વિધેયને
$\fn{EndSeq}$ કહે છે, પણ આગળ
$\fn{EndSequent}$ વાપરે છે; અહીં
બીજું નામ સતત રાખ્યું છે. આરંભિક
સિક્વન્ટવાળી બે-ઘટક ટપલ માટે
છેલ્લા નિયમનો દર્શાવેલો ઘટક નથી,
એથી એ વિધેયનું દર્શાવેલું સમીકરણ
માત્ર બિનઆરંભિક કિસ્સાને લાગુ પડે છે;
આરંભિક કિસ્સામાં~$0$ મૂકવાથી
પછીનું વિધાન સર્વત્ર અર્થપૂર્ણ બને છે.}
નીચેના વિધાનથી વ્યાખ્યાયિત
$\fn{Sequent}(s)$ ગુણધર્મ જુઓ:
\[
  \len{s} = 2 \land \bforall{i<\len{(s)_0} + \len{(s)_1}}{\fn{Sent}(((s)_0 \concat (s)_1)_i)}
\]
તે $s$ માટે ત્યારે અને માત્ર ત્યારે
સત્ય છે જ્યારે $s$ એ !!{sentence}sના
સિક્વન્ટની ગડેલ-સંખ્યા હોય. કેટલાક
કિસ્સા ઘણાં કઠિન છે. તેમને કેવી રીતે
સંભાળવા તેનો ઓછામાં ઓછો રૂપરેખાત્મક
ખ્યાલ આપીશું. હેતુ એ બતાવવાનો છે કે
``$\pi$ એ~$\Gamma$માંથી~$!A$ની
!!a{derivation} છે'' એ સંબંધ
$\pi$ અને~$!A$ની ગડેલ-સંખ્યાઓનો
આદિમ પુનરાવર્તી સંબંધ છે.
\end{explain}

\begin{prop}\ollabel{prop:followsby}
  ગડેલ-સંખ્યા~$p$ ધરાવતી
  !!{derivation}~$\pi$નું છેલ્લું અનુમાન
  યોગ્ય હોય ત્યારે અને માત્ર ત્યારે સત્ય
  થતો ગુણધર્મ~$\fn{Correct}(p)$ આદિમ
  પુનરાવર્તી છે.
\end{prop}

\begin{proof}
  $\Gamma \Sequent \Delta$ આરંભિક સિક્વન્ટ
  છે, જો કાં તો કોઈ !!a{sentence}~$!A$
  માટે $\Gamma \Sequent \Delta$ એ
  $!A \Sequent !A$ હોય,
  અથવા કોઈ પદ~$t$ માટે
  $\Gamma \Sequent \Delta$ એ
  $\emptyset \Sequent \eq[t][t]$ હોય.
  ગડેલ-સંખ્યાઓની ભાષામાં
  $\fn{InitSeq}(s)$ ત્યારે અને માત્ર ત્યારે
  સત્ય છે જ્યારે
  \begin{align*}
  \bexists{x < s}{} (\fn{Sent}(x) & \land
  s = \tuple{\tuple{x},\tuple{x}})
  \lor {}\\
  \bexists{t<s}{} (\fn{Term}(t) & \land
  s = \tuple{0, \tuple{\Gn{{\eq}(} \concat t \concat \Gn{,} \concat t \concat \Gn{)}}}).
  \end{align*}

  દરેક અનુમાન-નિયમ~$R$ માટે
  $\fn{FollowsBy}_R(p)$ સંબંધ આદિમ
  પુનરાવર્તી છે તે પણ બતાવવું પડે.
  આ $\fn{FollowsBy}_R(p)$ સંબંધ ત્યારે
  અને માત્ર ત્યારે સત્ય છે જ્યારે $p$ એ
  !!{derivation}~$\pi$ની ગડેલ-સંખ્યા
  હોય અને~$\pi$નું અંતિમ સિક્વન્ટ
  $\pi$ના સીધા ઉપ-!!{derivation}sમાંથી
  $R$ના યોગ્ય ઉપયોગ વડે ફલિત થતું હોય.

  એક સરળ કિસ્સો \RightR{\land} નિયમનો
  છે. $\pi$નો અંત યોગ્ય
  $\RightR{\land}$ અનુમાનથી થતો હોય,
  તો તે આ પ્રમાણે દેખાય:
  \begin{prooftree}
    \AxiomC{}
    \RightLabel{$\pi_1$}
    \Deduce$\Gamma \fCenter \Delta, !A$
    
    \AxiomC{}
    \RightLabel{$\pi_2$}
    \Deduce$\Gamma \fCenter \Delta, !B$

    \RightLabel{\RightR\land}
    \BinaryInf$\Gamma \fCenter \Delta, !A \land !B$
  \end{prooftree}
  એટલે, !!{derivation}~$\pi$નું છેલ્લું
  અનુમાન $\RightR{\land}$નો યોગ્ય ઉપયોગ
  ત્યારે અને માત્ર ત્યારે છે જ્યારે
  !!{sentence}sની શ્રેણીઓ $\Gamma$ અને
  $\Delta$, તેમજ બે !!{sentence}s
  $!A$ અને~$!B$ એવાં હોય કે
  $\pi_1$નું અંતિમ સિક્વન્ટ
  $\Gamma \Sequent \Delta, !A$,
  $\pi_2$નું અંતિમ સિક્વન્ટ
  $\Gamma \Sequent \Delta, !B$ અને
  $\pi$નું અંતિમ સિક્વન્ટ
  $\Gamma \Sequent \Delta, !A \land !B$ હોય.
  આને ગડેલ-સંખ્યાઓમાં વ્યક્ત કરીએ.
  $s = \Gn{\Gamma \Sequent \Delta}$
  હોય, તો $(s)_0 = \Gn{\Gamma}$
  અને $(s)_1 = \Gn{\Delta}$.
  તેથી $\fn{FollowsBy}_{\RightR{\land}}(p)$
  ત્યારે અને માત્ર ત્યારે સત્ય છે જ્યારે
\begin{align*}
& \bexists{g < p}{\bexists{d < p}{\bexists{a < p}{\bexists{b < p}{\quad}}}} \\
& \qquad \fn{EndSequent}(p) = 
  \tuple{g, d \concat 
    \tuple{\Gn{(} \concat a \concat \Gn{\land} \concat b \concat \Gn{)}}} \land {} \\
& \qquad \fn{EndSequent}((p)_1) = \tuple{g, d \concat \tuple{a}} \land {}\\
& \qquad \fn{EndSequent}((p)_2) = \tuple{g, d \concat \tuple{b}} \land {}\\
& \qquad (p)_0 = 2 \land \fn{LastRule}(p) = 10.
\end{align*}
અનુક્રમે આ પંક્તિઓ કહે છે કે
``ગડેલ-સંખ્યા~$g$ ધરાવતી
શ્રેણી~($\Gamma$), ગડેલ-સંખ્યા~$d$
ધરાવતી શ્રેણી~($\Delta$),
ગડેલ-સંખ્યા~$a$ ધરાવતું
!!a{formula}~($!A$) અને
ગડેલ-સંખ્યા~$b$ ધરાવતું
!!a{formula}~($!B$) અસ્તિત્વમાં છે'';
``$\pi$નું અંતિમ સિક્વન્ટ
$\Gamma \Sequent \Delta, !A \land !B$
છે''; ``$\pi_1$નું અંતિમ સિક્વન્ટ
$\Gamma \Sequent \Delta, !A$ છે'';
``$\pi_2$નું અંતિમ સિક્વન્ટ
$\Gamma \Sequent \Delta, !B$ છે'';
અને ``$\pi$નાં બે સીધાં
ઉપપ્રમાણો છે તથા છેલ્લો અનુમાન-નિયમ
$\RightR\land$ (નંબર~$10$) છે.''

$\pi$નું છેલ્લું અનુમાન
$\RightR\lexists$નો યોગ્ય ઉપયોગ હોય,
તો અને માત્ર તો જ !!a{formula}~$!A$,
ચલ~$x$, પદ~$t$ અને શ્રેણીઓ
$\Gamma$, $\Delta$ એવાં હોય કે
$\pi$નું અંતિમ સિક્વન્ટ
$\Gamma \Sequent \Delta, \lexists[x][!A]$
અને $\pi_1$નું અંતિમ સિક્વન્ટ
$\Gamma \Sequent \Delta,
\Subst{!A}{t}{x}$ હોય. તેથી
ગડેલ-સંખ્યાઓમાં
$\fn{FollowsBy}_{\RightR\lexists}(p)$
ત્યારે અને માત્ર ત્યારે સત્ય છે જ્યારે
\begin{align*}
  & \bexists{g<p}{\bexists{d<p}{\bexists{a<p}{\bexists{x<p}{\bexists{t<p}{\quad}}}}}\\
  & \qquad \fn{EndSequent}(p) = 
  \tuple{
    g,
    d \concat \tuple{\Gn{\lexists} \concat x \concat a}
  } \land {}\\
  & \qquad \fn{EndSequent}((p)_1) = 
  \tuple{
    g,
    d \concat \tuple{\fn{Subst}(a, t, x)}
  } \land {}\\
  & \qquad (p)_0 = 1 \land \fn{LastRule}(p) = 18.
\end{align*}
હવે $\fn{Correct}(p)$ને આ પ્રમાણે
વ્યાખ્યાયિત કરીએ:
\begin{multline*}
  \fn{Sequent}(\fn{EndSequent}(p)) \land {}\\ 
  [(\fn{LastRule}(p) = 1 \land 
  \fn{FollowsBy}_{\LeftR\Weakening}(p)) \lor \dots \lor {}\\
  (\fn{LastRule}(p) = 20 \land \fn{FollowsBy}_{\eq}(p)) \lor {}\\
  (p)_0 = 0 \land \fn{InitSeq}(\fn{EndSequent}(p))]
\end{multline*}
\sourcecorrection{OLINC-008}{સ્થિર અંગ્રેજી મૂળમાં
આરંભિક સિક્વન્ટના ગુણધર્મનું નામ
પહેલાં $\fn{InitSeq}$ છે, પણ આ
છેલ્લી પંક્તિમાં $\fn{InitialSeq}$ છે.
અહીં અગાઉ વ્યાખ્યાયિત નામ જ રાખ્યું છે.
આગળની સમજૂતીમાં અંતિમ સિક્વન્ટને
ભૂલથી~$d$નું કહેવામાં આવ્યું છે;
અહીં યોગ્ય ચલ~$p$ રાખ્યો છે.}
પહેલી પંક્તિથી~$p$નું અંતિમ સિક્વન્ટ
ખરેખર !!{sentence}sનું સિક્વન્ટ છે
તે નિશ્ચિત થાય છે. છેલ્લી પંક્તિ
$p$ માત્ર આરંભિક સિક્વન્ટ હોય તે
કિસ્સાને આવરી લે છે.
\end{proof}

\begin{prob}
  \olref[inc][art][plk]{prop:followsby}ની
  જેમ નીચેના ગુણધર્મો વ્યાખ્યાયિત કરો:
  \begin{enumerate}
  \item $\fn{FollowsBy}_{\Cut}(p)$,
  \item $\fn{FollowsBy}_{\LeftR\lif}(p)$,
  \item $\fn{FollowsBy}_{\eq}(p)$,
  \item $\fn{FollowsBy}_{\RightR{\lforall}}(p)$.
  \end{enumerate}
  છેલ્લા ગુણધર્મ માટે એ પણ બતાવવું પડશે
  કે ગડેલ-સંખ્યા~$p$ ધરાવતી
  !!{derivation}નું છેલ્લું અનુમાન
  વિશિષ્ટ ચલની શરત સંતોષે છે કે નહીં
  તે આદિમ પુનરાવર્તી રીતે ચકાસી શકાય.
  એટલે $\RightR{\lforall}$નો વિશિષ્ટ
  ચલ~$a$ અંતિમ સિક્વન્ટમાં આવતો નથી.
  \end{prob}
  

\begin{prop}
  \ollabel{prop:deriv}
  યોગ્ય !!{derivation}~$\pi$ની
  ગડેલ-સંખ્યા~$p$ હોય ત્યારે સત્ય
  થતો સંબંધ~$\fn{Deriv}(p)$ આદિમ
  પુનરાવર્તી છે.
\end{prop}

\begin{proof}
  !!^a{derivation}~$\pi$ ત્યારે યોગ્ય
  છે જ્યારે તેનું દરેક અનુમાન કોઈ
  નિયમનો યોગ્ય ઉપયોગ હોય, એટલે કે
  દરેક ઉપ-!!{derivation}નો અંત યોગ્ય
  અનુમાનથી થતો હોય. તેથી
  $\fn{Deriv}(p)$ ત્યારે અને માત્ર ત્યારે,
  જ્યારે
  \[
  \bforall{i<\len{\fn{SubtreeSeq}(p)}}{\fn{Correct}((\fn{SubtreeSeq}(p))_i)}.
  \]
  \sourcecorrection{OLINC-009}{સ્થિર અંગ્રેજી
  મૂળમાં આ અનુચ્છેદના ગદ્યમાં
  $\fn{Deriv}(d)$ લખ્યું છે, જ્યારે
  વિધાન અને સૂત્રનો ચલ~$p$ છે.
  મૂળ સૂત્રમાં $\fn{Correct}$ના ઘટકનું
  બંધ કૌંસ પણ છૂટી ગયું છે; અહીં
  બંને અસંગતતાઓ સુધારી છે.}
\end{proof}


\begin{prop}
$\Gamma$ એ !!{sentence}sનો આદિમ
પુનરાવર્તી ગણ હોય એમ ધારો. ત્યારે
$\Prf[\Gamma](x, y)$ સંબંધ આદિમ
પુનરાવર્તી છે. તેનો અર્થ એ છે કે
$x$ એ કોઈ સાન્ત
$\Gamma_0 \subseteq \Gamma$ માટે
$\Gamma_0 \Sequent !A$ની
!!a{derivation}~$\pi$નો સંકેત છે અને
$y$ એ~$!A$ની ગડેલ-સંખ્યા છે.
\end{prop}

\begin{proof}
``$y \in \Gamma$'' આદિમ પુનરાવર્તી
વિધાન~$R_\Gamma(y)$ વડે આપવામાં
આવ્યું હોય એમ ધારો. આપણે બતાવવું છે
કે $\Prf[\Gamma](x, y)$ આદિમ
પુનરાવર્તી છે. આ સંબંધ ત્યારે અને
માત્ર ત્યારે સત્ય છે જ્યારે $y$ એ
વાક્ય~$!A$ની ગડેલ-સંખ્યા
હોય અને $x$ એ અંતિમ સિક્વન્ટ
$\Gamma_0 \Sequent !A$ ધરાવતી
$\Log{LK}$-!!{derivation}નો સંકેત હોય.

અગાઉના વિધાન પ્રમાણે, $x$ એ
$\Log{LK}$માં યોગ્ય
!!{derivation}~$\pi$નો સંકેત હોય
ત્યારે સત્ય થતો
$\fn{Deriv}(x)$ ગુણધર્મ આદિમ
પુનરાવર્તી છે. $x$ એવો સંકેત હોય,
તો $\fn{EndSequent}(x)$ એ~$\pi$ના
અંતિમ સિક્વન્ટનો સંકેત છે. તેથી
$(\fn{EndSequent}(x))_0$ અંતિમ
સિક્વન્ટની ડાબી બાજુનો સંકેત છે
અને $(\fn{EndSequent}(x))_1$
જમણી બાજુનો સંકેત છે. આમ,
``$\pi$ના અંતિમ સિક્વન્ટની જમણી
બાજુ માત્ર~$!A$ છે'' તે
$\len{(\fn{EndSequent}(x))_1} = 1
\land ((\fn{EndSequent}(x))_1)_0 = y$
વડે વ્યક્ત થાય છે.
\sourcecorrection{OLINC-010}{સ્થિર અંગ્રેજી મૂળના
આ ગદ્ય-સૂત્રમાં છેલ્લે $y$ને બદલે
$x$ છે; પરંતુ તરત નીચેની
$\Prf[\Gamma](x,y)$ની વ્યાખ્યામાં
યોગ્ય રીતે $y$ છે. અહીં ગદ્ય-સૂત્રને
એ જ વ્યાખ્યા સાથે મેળમાં કર્યું છે.}
$\pi$ના અંતિમ સિક્વન્ટની ડાબી
બાજુ તો આપોઆપ સાન્ત છે; માત્ર
તેનું દરેક વાક્ય~$\Gamma$માં છે
તે વ્યક્ત કરવું પડે. આમ
$\Prf[\Gamma](x, y)$ને આ પ્રમાણે
વ્યાખ્યાયિત કરી શકીએ:
\begin{align*}
\Prf[\Gamma](x, y) \defiff {}&
\fn{Deriv}(x) \land {} \\
& \bforall{i <
  \len{(\fn{EndSequent}(x))_0}}{R_\Gamma(((\fn{EndSequent}(x))_0)_i)} \land {}\\
& \len{(\fn{EndSequent}(x))_1} = 1 \land ((\fn{EndSequent}(x))_1)_0 = y.
\end{align*}
\end{proof}

\end{document}
